\sscp{Stresemann (1927)}{07-01-01}
In one of the earliest systematic attempts at subgrouping,
\citet{Stresemann1927} identified a number of mergers
and losses to define his “Proto-Ambon” (“Ur-Ambon”) subgroup,
and used those to exclude other languages on the periphery from this subgroup.
Within Proto-Ambon \citeauthor{Stresemann1927} grouped “Sub-Buru” closer to “Sub-Ambon”,
and further from “Sub-Seran”.
However, he also noted (p. 7) that of the three innovations
he identified for defining “Sub-Ambon” (merger of *d/*R/*l >
**l; common loss of *k; presence of support vowel \it{-e} after
nouns ending in a final consonant),
“none have validity” for Sub-Buru languages.
Kayeli (“Kajeli”), in eastern Buru, was seen as bridging between
Sub-Buru and Sub-Ambon, sharing features with both groups (p. 8).
Sub-Seran was undefined, or negatively defined as Ambon languages
that did not fit in the other two groups (p. 7).
\citeauthor{Stresemann1927}’s subgroups are shown in \frf{07-06}.

\begin{figure}[h]\centering
	\includegraphics[width=\columnwidth]{StresemannMap}
	\caption[Location of Stresemann’s (1927) subgroups in central Maluku]
	{Location of \citeauthor{Stresemann1927}’s (\citeyear{Stresemann1927}) Maluku subgroups}
	\label{fig:07-06}
\end{figure}

Since we separate \tsc{Sula-Buru} from \tsc{Ambon-Seram},
but \citeauthor{Stresemann1927} combines Sub-Buru with Sub-Ambon
and Sub-Seran together under his Proto-Ambon,
it begs the question about what we interpret differently from
the innovations \citeauthor{Stresemann1927} identified that define his higher
level Proto-Ambon.\footnote{
		\citeauthor{Stresemann1927}’s Sub-Buru included Buru
		and Ambelau, but not Sanana
		and Mangole north of Buru.
		Our \tsc{Sula-Buru} subgroup includes Ambelau,
		Buru, Lisela, Sanana and Mangole.
		Our \tsc{Ambon-Seram} subgroup is roughly equivalent to \citeauthor{Stresemann1927}'s
		Sub-Ambon combined with his Sub-Seran,
		but excluding his Sub-Buru.}						
To properly unravel his arguments
one must interpret his symbols
and proto-forms according to MP historical phonology
and reconstructions as they were understood at the time he was
writing.\footnote{
		For example, \citeauthor{Stresemann1927}'s *g{\Q} = PMP *j (he gives the etyma for *ŋajan ‘name’,
		*pajay ‘rice’, *pija ‘how many’, *ləjaw ‘sun’),
		his *r = PMP *R,
		his *d{\Q} = PMP *z (*zalan ‘path’, *quzan ‘rain’), etc.
		But \citeauthor{Stresemann1927}'s *ḍ is also sometimes PMP *R (*tabuRi ‘trumpet’),
		his *d is usually PMP *d,
		but sometimes PMP *j (*uləj ‘worm’) or *z (*qazay ‘jaw,
		chin’), and so forth.}
Once his symbols
and reconstructions are updated according to a modern understanding
of MP phonology, many of his arguments do not stand up to scrutiny.
Furthermore, many of his proposed mergers
and losses are now known to be found more widely
and are not exclusive innovations for Proto-Ambon.
\citet[14--102]{Stresemann1927} identifies the following features
as defining Proto-Ambon, with his summary of sound correspondences,
mergers and loss of consonants given on his pp. 88--89.
Supporting data for the discussion below is found in the following
sections and chapters.

\begin{itemize}

\item PMP *l/*j > Proto-Ambon **l; PMP *R = **r; PMP *d/*z > **d.
But \citet[17]{Stresemann1927} notes that these have different
correspondences for Sub-Buru, and that *j is problematic for Sub-Buru.
Sub-Ambon also merges PMP *R/*l/*j > **l,
whereas our \tsc{Sula-Buru} data show that reflexes of PMP *d,
*z, *R, *l, and *j must be kept distinct.
The merger of PMP *d/*z > Proto-Ambon **d for Sub-Buru shown on p.
88 is also not well supported by his own data on pp.
16--18.\footnote{
		For example, he correctly listed Buru (``Tifu'')
		\it{aa-} ‘jaw, chin’ < Proto-Ambon **ada (PMP *qazay),
		which we recognise as showing PMP *z > Buru {\0} -- the correct
		sound correspondence for Buru,
		whereas PMP *d > Buru \it{r}.
		Instead of using the symbol **d{\Q},
		he used **d for **ada ‘jaw,
		chin’, and sometimes for *zalan ‘path’ > **dalan;
		*quzan ‘rain’ > **udan,
		making it harder for him to see that \tsc{Sula-Buru} does not
		actually merge PMP *d
		and *z.}
\item PMP *k/*g > Proto-Ambon **k (p.
88) is incorrect for \tsc{Sula-Buru}.
Buru retains separate reflexes,
as seen in \citeauthor{Stresemann1927}'s own data on p. 68.
Further evidence can be found in \citet{GrimesGrimes2020}.

\item PMP *nd/*nz/*nj/*nt/*ns > Proto-Ambon **nd is problematic for
\tsc{Sula-Buru} in many ways,
including the scarcity of reconstructed etyma
and reflexes that support them,
which leads \citeauthor{Stresemann1927} to appeal to data from Kayeli to support
the merger for Sub-Buru.
We place Kayeli within \tsc{Ambon-Seram}.
The scarce data available do not support this merger for
\tsc{Sula-Buru}.\footnote{
		For Buru, the scarce data available suggest that *nd > \it{d},
		*nt > \it{t}, *ns > \it{s},
		showing a consistent loss of the nasal
		and preservation of the following consonant.
		Normally for Buru, *d > \it{r},
		*t = \it{t}, *s = \it{s}.
		Ambelau also shows *nt > \it{t} in at least two items.}

\item *ŋk/*ŋg > Proto-Ambon **ŋg has too few etyma
and even fewer reflexes to be able to verify actual patterns.
Some of the evidence provided is based on Malay loans,
such as \it{tunggu} [tuŋgu] ‘wait’ (p. 67).

\item *-aw/*-ay > Proto-Ambon **a is correct for \tsc{Sula-Buru},
but not for \tsc{Ambon-Seram} (see \srf{10-02}).

\item *-uy > Proto-Ambon **u is slightly problematic for both \tsc{Sula-Buru}
and \tsc{Ambon-Seram}, as traces of the full vowel-glide sequence
must be preserved to account for some data in both subgroups.

\item final *p{\#}/*k{\#} > Proto-Ambon **ʔ{\#} in nouns has no evidence
to support it from \tsc{Sula-Buru}.
\citeauthor{Stresemann1927} has not recognised the productive nominal suffix \it{-t},
and concludes that Proto-Ambon **ʔ > \it{t}{\#} in Buru. (p.
75). The apparent retention of [ʔ] in \citeauthor{Stresemann1927}’s Ambelau
and Kayeli data are probably misleading transcription devices
from \citeauthor{Stresemann1927}’s sources, since none of our modern sources have glottals for those same
words in either Ambelau or Kayeli.
Otherwise, our Kayeli and Ambelau data (collected in the late
1980s) are the same as his.

\item *mb/*mp > Proto-Ambon **mb is correct for \tsc{Sula-Buru}.
\tsc{Ambon-Seram} is better served with *mb/*mp > **mp,
reflecting the devoicing that is found from that cluster,
even in languages that retain voiced PMP *b = \it{b}.
Furthermore, \citet[26]{Blust1981} sees *mb/*mp > **mb as an
innovation spanning the wider CMP region,
so it has little subgrouping value for Proto-Ambon.

\item *n/*ñ > Proto-Ambon **n is so widespread as to have no real subgrouping value.
And there are also trace retentions of *ñ in raised vowels in
some languages that suggest *ñ must still be retained for some
etyma.\footnote{
		\citet[12]{Collins1983} misinterpreted
		\citeauthor{Stresemann1927}'s \citeyear{Stresemann1927}
		symbols on page 88 as a merger of PMP *ŋ/*n, rather than *n/*ñ.
		\citeauthor{Stresemann1927} gives the correspondences for PMP *ŋ separately on page 89.
		The error is an easy one to make,
		since \citeauthor{Stresemann1927}'s symbols for both *ŋ
		and *ñ involve two similar diacritics above ‹n› in a font
		that makes it hard to see the diacritics.
		He represents *ŋ as ‹ṅ› with a dot above \it{n},
		and he represents *ñ as ‹ń› with an acute accent above \it{n}.}

\item *h > Proto-Ambon {\0} is so widespread as to have no subgrouping
value at this lower level.

\item Loss of final consonants in verbs: not an exclusive innovation
to Proto-Ambon, so it has no subgrouping value.

\item Loss of final consonants in nouns,
except *R, *l, and *t.
It turns out that other final consonants are also retained,
the retentions are found in the same etyma beyond Proto-Ambon
thus not being exclusive,
and \tsc{Sula-Buru} and \tsc{Ambon-Seram} have different patterns
regarding the retention of final *R,
all of which are shown in \srf{07-03}.
\end{itemize}

Careful examination of \citeauthor{Stresemann1927}’s features defining Proto-Ambon
leaves no innovations that are shared exclusively between
Sub-Buru and Sub-Ambon (with or without Sub-Seran).
We already noted that \citet[7]{Stresemann1927} acknowledged
that none of the innovations defining Sub-Ambon are found in Sub-Buru.
As we further show in \srf{07-04}
and chapters \ref{ch:SB} and \ref{ch:AS}, the historical phonology
and morphosyntax show \tsc{Sula-Buru}
and \tsc{Ambon-Seram} are two distinct subgroups,
in which their differences are far more striking than their similarities.

\citeauthor{Stresemann1927} excluded from Proto-Ambon the languages of the Sula
archipelago in the northwest on the basis of grammar.
In the east he excluded “Proto-Gorom” (“Ur-Goram” -- which included
the “Seranlaut”\footnote{
		\citeauthor{Stresemann1927}’s “Seranlaut” is mostly from
		mainland varieties of the Geser-Gorom language,
		and seems to include some others as well.}
languages of east
Seram, Geser-Gorom and Watubela), Banda (“Eli-Elat”),
and Masiwang (“Bonfia”) the latter of which was seen as sharing
some features with Proto-Ambon,
and diverging in others. See \frf{07-06}.
Our \tsc{Seram-Tanimbar-Bomberai} subgroup includes
and extends beyond \citeauthor{Stresemann1927}’s Proto-Gorom,
and includes several languages he excluded from his Proto-Ambon
subgroup, such as Masiwang (“Bonfia”).
